% !TEX root = ../main.tex
\chapter{Conclusion and Future Work}
\label{ch:conclusion}

Chapter~\ref{ch:discussion} interpreted the findings and judged whether the study met its objectives. The tables of Chapter~\ref{ch:results} form the record.

\dissertationSubheading[sec:summary-findings]{Answers to the research questions}

The thesis proposed EndorseRank, PageRank with the edge weights of Adaptive Weighted PageRank (AWP) on a latest-allowance endorsement graph, and evaluated it beside AWP itself\cend{3}, a coupled walk over both edge types and the raw degrees these walks smooth, on GMX~V2 wallets and spenders on Arbitrum One from December~2025 to May~2026, with a registered replication on June to August~2026.

RQ1 asked whether EndorseRank and AWP rank the wallets of one sample differently. They do, and H1 is supported, largely by construction (Section~\ref{sec:rank-divergence}). Much of the difference, and of the agreement, comes from EndorseRank's tie blocks, so the claim that the two reflect different kinds of reputation rests mainly on the holdouts of RQ4.

RQ2 asked how closely each score agrees with allowance and transfer statistics in the same window. Each score follows the edges it walks on (Table~\ref{tab:alignment-hybrid}), which supports H2 only narrowly, because the proxies come from the same events as the graphs; since EndorseRank's coefficients are set by its ties, no conclusion is drawn from the difference between its allowance and transfer coefficients.

RQ3 asked what each score costs to compute. The cost follows the edge count, so EndorseRank, on the sparsest graph, is the cheapest score, and the coupled walks cost more than AWP (Table~\ref{tab:benchmark-runtime}); H3 is supported, largely by construction. The pool stages could not test edge growth, because the second-tier extraction was not completed (Table~\ref{tab:benchmark-scaling}).

RQ4 asked whether scores frozen at $t_1$ predict new approvals and new transfer senders on the spenders at $t_2$, and whether any PageRank predicts them better than the raw degree it smooths. The first part holds, and the edge type decides which label a score anticipates (Table~\ref{tab:holdout-spenders}). The second does not: the raw in-approve degree predicts new approvals, and the raw transfer in-degree new senders, better than every walk in both windows (Tables~\ref{tab:holdout-diff} and~\ref{tab:fresh-posthoc}). For these labels, on this depth-one sample, propagation adds no predictive value over the local count (Section~\ref{sub:theory-closed-form}).

RQ5 asked whether a coupled PageRank over both edge types improves on a walk over the transfers alone, and whether it ranks wallets better than either layer walked alone. Both answers are no, each under a rule fixed in advance. Under the rule of 14~September, C-PR fails both conditions of the primary part and the allowance condition of the secondary part (Section~\ref{sec:coupled-results}; Table~\ref{tab:holdout-diff}). Its exploratory gain over the transfer walk on new approvals in the spring replicated in the registered June--August window, but a loss on new senders larger than the registered margin could not be ruled out, so the first answer is no; against AWP the loss on new senders was clear (Table~\ref{tab:fresh-contrasts}). The seeded ablation matches the transfer walk: allowance edges contribute when the walk moves along them, and not as a seed.

In a pre-specified post hoc analysis that adds no research question (Section~\ref{sub:finding-neutral}), the approval count at the freeze ranked traders above the transfer count in both windows on later liquidation avoidance, the outcome built from neither edge type that bears on counterparty risk ($\Delta\tau=+0.150$ in June--August; Section~\ref{sub:neutral-label-results}). This lead holds at count level only: the like-for-like PageRank pair is not resolved, account type is an unresolved confound because most approval-receiving traders are contract accounts, and the signal is lower on the profitable share of closes, so it ranks forced-liquidation avoidance, not trading success.

Claims~C1--C6 (Section~\ref{sec:empirical-claims}) state these findings formally.

\dissertationSubheading[sec:contribution-knowledge]{Original contribution to knowledge}

In the terms of the literature reviewed in Chapter~\ref{ch:literature}, the five contributions of Section~\ref{sec:contributions} establish the following.

Before this work, graph reputation on public ledgers ranked wallets by their transfers \cend{3,4}, by how risk spreads through transaction networks \cend{12}, or by explicit ratings and stakes \cend{25,26}, and the reputation tooling deployed in DeFi scored the holder of an address from credentials \cend{64}. None of it used the ERC-20 allowance, the on-chain record of a deliberate delegation of spending power \cend{18}, as the edge of a reputation ranking; approvals had been used to cluster the addresses of one entity \cend{40} and studied as a security risk \cend{73}. None evaluated on-chain reputation scores on one cohort against the checks of each construct and against the raw degrees they smooth, with confidence intervals, difference tests and a temporal holdout. This thesis supplies both; PageRank and AWP's weighting are taken from the literature \cend{3,15}, EndorseRank changes only the edge, and the coupled operator changes the walk.

The new knowledge has four parts. What an on-chain reputation score measures depends on the edges it uses, and this can be verified on one cohort with quantified uncertainty. The allowance edge yields a score that is cheaper to update, differs from the transfer-based score, and ranks later endorsements of spenders in two windows, a label on which the scores built on transfers alone stay low. Post hoc and at count level, the approval count also ranked traders' later avoidance of forced liquidation, though not their trading success, above the transfer count; the PageRank pair is not resolved, and account type may explain the lead. Two further steps that the literature suggested did not improve results on these data: propagation predicts the future growth of the graph no better than the raw degrees it smooths, and a weighted combination of the two edge types beats neither its better single layer nor, under the registered rule, the transfer-only walk. Finally, the first of these nulls has an analytic explanation for EndorseRank: on the depth-one allowance graphs of this sample it is, up to a constant, a linear function of the owner-normalized in-strength, which divides each owner's endorsement over all the spenders it approves (Section~\ref{sub:theory-closed-form}); multi-hop propagation remains untested. The construct-validity framing \cend{22,33} that makes these claims testable is part of the contribution, as is the open pipeline whose published summaries let others check every table, since whether decentralized systems miscalculate credit scores is a matter of public interest \cend{2}.

The scope is deliberately narrow. The thesis does not show that on-chain reputation can replace collateral \cend{1,9}, that the results apply beyond Arbitrum One and the window studied, or that any score measures creditworthiness, default risk or trading skill, and it does not claim that EndorseRank resists Sybils: under uniform teleportation a closed farm of attacker addresses can inflate it cheaply (Section~\ref{sec:sybil-cost-model}).

\dissertationSubheading[sec:future-work]{Future work}

Each suggestion below corresponds to a limitation identified in Chapter~\ref{ch:discussion}. Any new score or weighting it leads to should be fixed before it is scored, replicated on registered new data and checked out of window against degree baselines.

\begin{itemize}
\item A registered test of the liquidation-avoidance lead. Claim~C6 is post hoc and should be tested on a later window whose logs have not yet been extracted, with the count pair (in-approve degree against transfer in-degree) and one PageRank pair fixed in advance: AWP against a walk over AWP's transfers plus an allowance layer with AWP's weights and restarts, so that only the edge type changes. Account type, from a code lookup of every trader, would enter as a baseline, and the cohort would be sized to detect a lead half as large as the post hoc one.
\item Propagation on a deeper sample. On the one-hop ego networks studied, EndorseRank reduces to a reweighted approval count (Section~\ref{sub:theory-closed-form}); a sample extracted two or more hops out would test whether any propagated score adds out-of-window value over the local count, on arrival labels and on labels that are not counts of arrivals.
\item Complete the second-tier extraction. The expanded address pool has so far only varied the number of nodes; with the logs of its supplemental addresses extracted, we could measure how running time grows with the number of edges and turn the bounded claim for RQ3 into a general scaling claim.
\item Validate the results in other contexts. Running the same analysis on other rollups and on Layer-1 networks would show whether the construct checks and holdout results extend beyond Arbitrum One; because DeFi is heterogeneous \cend{1,8}, this is the main threat to external validity.
\item Extend the time window and use a credit-linked holdout. A longer window would show whether the Kendall $\tau$ coefficients stay stable across market conditions, and a credit-linked holdout on the owners, feasible where lending protocols cover enough of the cohort, whether either graph carries credit-risk information.
\item Value weights. With $b=1$ on raw base units, AWP's value weight counts every approval or transfer of a few base units or more about once. A value of $b$ set on price-normalized amounts would let size matter again, and the spring holdout suggests that it might: the allowance layer weighted by raw amounts ranked both spender labels better than EndorseRank.
\item Adversarial evaluation. The Sybil cost model is analytical, backed by computations with the thesis's own PageRank code on the observed EndorseRank graph, and no attack was carried out on chain (Section~\ref{sec:sybil-cost-model}). Restricting teleportation and dangling redistribution to a vetted seed set, as TrustRank does \cend{17}, would remove the uniform inflow that closed farms exploit and should be evaluated under the same attacks. Empirical red-teaming with wash approvals, coordinated spender farms and strategic revocation would measure the realized cost-effectiveness of attacks under live fee markets and detection heuristics \cend{11}, and would test whether the Sybil-adjusted proxies remain informative under active adversaries.
\item Coupling beyond a fixed transition-level mixture. Only a fixed $\lambda$ with uniform teleportation on raw-amount layers was tried. Node-dependent or learned weights, per-layer damping, mixtures of the two stationary vectors, and layers weighted as in EndorseRank and AWP remain untested; any such design should report construct-specific checks and keep the auditability that sets graph propagation apart from opaque learned scores \cend{2,14}.
\item Richer authorization semantics. Where ERC-20 approvals are sparse, adding protocol-specific delegation logs to the edge definition may improve coverage.
\item Composition with attestation tooling. Filtering addresses through a Passport-style gate \cend{64}, which answers whether an address belongs to a distinct person, before EndorseRank orders them by the authority others delegate to them is a plausible deployment; it should be evaluated on labeled Sybil data, with the gate's own weaknesses \cend{65,66} in view.
\item Deployment studies. Integrating EndorseRank with risk oracles and governance tooling would test whether its runtime profile can sustain production refresh cadences. Controlled experiments on fee tiers or monitoring cadence would let evidence replace the adoption hypotheses of Section~\ref{sub:deploy-industry}.
\end{itemize}

\dissertationSubheading[sec:closing-reflection]{Closing reflection}

In adversarial, pseudonymous markets on-chain reputation does not remove the need for collateral \cend{9}, but it can make behavioral structure readable, and which structure it reads depends on the edges used. A transfer records economic flow and an allowance records delegated authority. PageRank will rank either graph, yet one ranking does not serve both constructs, and a walk over both graphs did no better on either than the right single graph. For a protocol or researcher who needs a fast, auditable signal of which spenders owners trust, the allowance edge is the one to use on data like those studied here. The approval count is its cheapest form, and EndorseRank is the propagated form for those who want the endorsers to count as well, with its exposure to cheap inflation by closed farms in mind (Section~\ref{sec:sybil-cost-model}). For describing transfer hubs and anticipating their inflow, the transfer edge is the one to use. For screening counterparties against forced liquidation, the approval count ranked traders in these data and the transfer count, which ranked how much traders trade, did not; this post hoc, count-level result concerns liquidation avoidance, not trading success, is unresolved at the PageRank level, may reflect account type, and awaits a registered test.
